\section*{अध्याय \ref{ch:b3:rate-theories} --- अभिक्रिया दरों के सिद्धांत}

\begin{solution}{exo:b3:rate-theories:1}
$\bar v_{\mathrm{rel}} = \sqrt{8 \times 1.381\times10^{-23} \times 298/(\pi \times 18.46 \times 1.661\times10^{-27})} = \qty{585}{m/s}$।
$\sigma\bar v_{\mathrm{rel}}N_A = 0.43\times10^{-18} \times 585 \times 6.022\times10^{23} = \qty{1.51e8}{m^3.mol^{-1}.s^{-1}} =
\qty{1.5e11}{L.mol^{-1}.s^{-1}}$। मापित: $2.07\times10^{-12} \times 6.022\times10^{23}\times10^{-3} =
\qty{1.25e9}{L.mol^{-1}.s^{-1}}$, अतः $P = 0.008$: सौ में एक से भी कम संघट्ट का अभिविन्यास
सही होता है (NO के O को \ce{O3} के किसी अंतस्थ O से मिलना होता है)।
\end{solution}

\begin{solution}{exo:b3:rate-theories:2}
विलयन में $\Delta^{\ddagger}H^\circ = 75.0 - 2.48 = \qty{72.5}{kJ/mol}$; द्विआण्विक गैस अभिक्रिया के लिए
$75.0 - 4.96 = \qty{70.0}{kJ/mol}$।
\end{solution}

\begin{solution}{exo:b3:rate-theories:3}
डील्स--ऐल्डर: प्रबल रूप से ऋणात्मक (दो अणु एक चक्रीय संकुल बनाते हैं)।
\ce{N2} की हानि: धनात्मक (एक आबंध ढीला होता है, स्वातंत्र्य मिलता है)। ध्रुवीय विलायक में
बनते आयन: ऋणात्मक, विकसित होते आवेश अपने चारों ओर विलायक को क्रमित करते हैं।
\end{solution}

\begin{solution}{exo:b3:rate-theories:4}
जल: $8 \times 8.314 \times 298.15/(3 \times \num{0.890e-3}) = \qty{7.4e9}{L.mol^{-1}.s^{-1}}$; हेक्सेन
\qty{2.2e10}{L.mol^{-1}.s^{-1}}। हेक्सेन, जो तीन गुना कम श्यान है: अणु तीन
गुना तेज़ी से विसरित होते हैं।
\end{solution}

\begin{solution}{exo:b3:rate-theories:5}
$\ln(k_2T_1/k_1T_2) = \ln(10 \times 298.15/318.15) = 2.238$; $\Delta^{\ddagger}H^\circ = 8.3145 \times 2.238/(1/298.15 -
1/318.15) = \qty{88.2}{kJ/mol}$। $\Delta^{\ddagger}S^\circ = R[\ln(kh/k_BT) + \Delta^{\ddagger}H^\circ/RT] = 8.3145 \times
(-36.365 + 35.596) = \qty{-6.4}{J.K^{-1}.mol^{-1}}$ (\qty{298.15}{K} पर)। $E_a = R\ln10/(1/298.15 - 1/318.15) =
\qty{90.8}{kJ/mol} \approx \Delta^{\ddagger}H^\circ + RT$, माध्य ताप पर।
\end{solution}

\begin{solution}{exo:b3:rate-theories:6}
$\mu_{\mathrm{OH}} = 16/17$, $\mu_{\mathrm{OD}} = 32/18$, तरंग-संख्याओं का अनुपात $\sqrt{1.889} = 1.374$:
$\tilde\nu_{\mathrm{OD}} = \qty{2661}{cm^{-1}}$। $k_{\mathrm H}/k_{\mathrm D} = \exp(1.4388 \times 996/(2 \times 298)) = \eu^{2.40} = 11$।
\end{solution}

\begin{solution}{exo:b3:rate-theories:7}
अधिकतम $\exp(1.4388 \times 808/(2 \times 298)) = \eu^{1.95} = 7.0$। 40 का मान
\omterm{def:b3:quantum-model-systems:oscillator}{शून्य-बिंदु ऊर्जाओं} से नहीं आ सकता: हाइड्रोजन अवरोध से \omterm{def:b3:quantum-model-systems:tunnelling}{सुरंगन} करता है, जो ड्यूटेरियम,
दुगुना भारी होने से, कहीं कम करता है।
\end{solution}

\begin{solution}{exo:b3:rate-theories:8}
$\ce{S2O8^{2-}}$ और \ce{I-}: $z_{\mathrm A}z_{\mathrm B} = +2$, $\log(k/k_0) = 2 \times 0.51 \times 2 \times 0.10 = 0.20$: $k$,
1.6 से गुणा हो जाता है। एस्टर और \ce{OH-}: एक साथी उदासीन, कोई प्रभाव नहीं।
$\ce{[Co(NH3)5Br]^{2+}}$ और \ce{OH-}: $z_{\mathrm A}z_{\mathrm B} = -2$, $k$, $10^{-0.20} = 0.63$ से गुणा हो जाता है।
\end{solution}

\begin{solution}{exo:b3:rate-theories:9}
आरंभिक, प्रवेश घाटी में (अभिकारक-सदृश संक्रमण अवस्था, हैमंड
अभिगृहीत के अनुसार)। उत्क्रम अभिक्रिया, जो ऊष्माशोषी है, के अपने
पथ पर पश्च अवरोध है: H--F आबंध में कंपनिक ऊर्जा उसे सबसे अधिक बढ़ावा देती है।
\end{solution}

\begin{solution}{exo:b3:rate-theories:10}
$\Delta^{\ddagger}H^\circ = -Rb = \qty{61.9}{kJ/mol}$, मानक त्रुटि $R \times 31 = \qty{0.26}{kJ/mol}$ और
95\,\% अंतराल $\pm\qty{0.8}{kJ/mol}$ के साथ।
\[
  \Delta^{\ddagger}S^\circ = R(16.528 - 23.760) = \qty{-60.1}{J.K^{-1}.mol^{-1}},
\]
मानक त्रुटि $R \times 0.103 = \qty{0.86}{J.K^{-1}.mol^{-1}}$ और 95\,\% अंतराल $\pm\qty{2.7}{J.K^{-1}.mol^{-1}}$ के साथ।
\end{solution}

\begin{solution}{exo:b3:rate-theories:11}
$\ln k = \ln(k_BT/h) + \Delta^{\ddagger}S^\circ/R - \Delta^{\ddagger}H^\circ/RT$ से,
\[
  \frac{\dd\ln k}{\dd T} = \frac1T + \frac{\Delta^{\ddagger}H^\circ}{RT^2}, \qquad
  E_a = RT^2\frac{\dd\ln k}{\dd T} = \Delta^{\ddagger}H^\circ + RT .
\] तब $A = k\eu^{E_a/RT} = (k_BT/h)\eu^{\Delta^{\ddagger}S^\circ/R}\eu^{-\Delta^{\ddagger}H^\circ/RT}\eu^{(\Delta^{\ddagger}H^\circ + RT)/RT} =
\eu\,(k_BT/h)\eu^{\Delta^{\ddagger}S^\circ/R}$। $\Delta^{\ddagger}S^\circ = 0$ के साथ: $A = 2.718 \times \num{6.21e12} = \qty{1.7e13}{s^{-1}}$, जो
प्रेक्षित गुणकों की कोटि है; इससे ऊपर के मानों का अर्थ है धनात्मक $\Delta^{\ddagger}S^\circ$, एक ढीली
संक्रमण अवस्था।
\end{solution}

\begin{solution}{exo:b3:rate-theories:12}
प्रति एकक आयतन \omterm{def:b3:partition-functions:partition-function}{आण्विक विभाजन फलनों} के साथ: $q^{\mathrm T}/V = (2\pi mk_BT)^{3/2}/h^3$, A, B और संकुल (द्रव्यमान $m_{\mathrm A} + m_{\mathrm B}$) के लिए,
और संकुल के लिए $q^{\mathrm R} = 8\pi^2Ik_BT/h^2$
($\sigma = 1$)। स्थानांतरणीय अनुपात $h^3/(2\pi\mu k_BT)^{3/2}$ है, अतः
\[
  k = \frac{k_BT}{h}\cdot\frac{h^3}{(2\pi\mu k_BT)^{3/2}}\cdot\frac{8\pi^2\mu d^2k_BT}{h^2}\,\eu^{-E_0/RT}
    = \frac{8\pi^2d^2(k_BT)^{1/2}}{(2\pi)^{3/2}\mu^{1/2}}\eu^{-E_0/RT}
    = \pi d^2\sqrt{\frac{8k_BT}{\pi\mu}}\,\eu^{-E_0/RT},
\]
प्रति अणु-युग्म; $N_A$ से गुणा करने पर \cref{thm:b3:rate-theories:collision} मिलता है,
$\sigma = \pi d^2$ के साथ।
\end{solution}

\begin{solution}{pb:b3:rate-theories:1}\emergencystretch=3em
\textbf{1.} $\mu_{\mathrm{CH}} = 12 \times 1.00783/13.00783 = \qty{0.9297}{u}$; $\mu_{\mathrm{CD}} = 12 \times 2.01410/14.01410 =
\qty{1.7246}{u}$।
\textbf{2.} $2917/\sqrt{1.7246/0.9297} = 2917/1.3620 = \qty{2142}{cm^{-1}}$, मापे गए 2109 से 1.5\,\%
ऊपर (अनावर्तिता और \ce{CD4} के चारों आबंधों का युग्मन)।
\textbf{3.} $\frac12 \times 2917 = \qty{1458.5}{cm^{-1}}$ और $\frac12 \times 2109 = \qty{1054.5}{cm^{-1}}$; अंतर
\qty{404}{cm^{-1}} $= \qty{4.83}{kJ/mol}$।
\textbf{4.} वह अभिक्रिया निर्देशांक बन जाता है: उसकी \omterm{def:b3:quantum-model-systems:oscillator}{शून्य-बिंदु ऊर्जा} लुप्त हो जाती है।
\textbf{5.} D के लिए अवरोध \qty{4.83}{kJ/mol} ऊँचा है।
\textbf{6.} $\exp(1.43878 \times 404/310.15) = \eu^{1.874} = 6.5$।
\textbf{7.} $\exp(1.43878 \times 404/298.15) = 7.0$।
\textbf{8.} वास्तविक संक्रमण अवस्था में तनन की \omterm{def:b3:quantum-model-systems:oscillator}{शून्य-बिंदु ऊर्जा} का कुछ भाग
अन्य विधाओं में बचा रहता है, जो प्रभाव को घटाता है; केवल \omterm{def:b3:quantum-model-systems:oscillator}{शून्य-बिंदु ऊर्जाएँ}
इस मान से अधिक नहीं हो सकतीं (\omterm{def:b3:quantum-model-systems:tunnelling}{सुरंगन} हो सकता है)।
\textbf{9.} बहुत कम: वे द्वितीयक प्रभाव हैं, अधिक से अधिक प्रति ड्यूटेरियम लगभग 10\,\% की
कोटि के।
\textbf{10.} कि C--H आबंध दर-निर्धारक चरण में टूटता है, और
तनन पूर्णतः अभिक्रिया निर्देशांक में बदल जाता है।
\textbf{11.} प्रथम-कोटि विलोपन: $t_{1/2} = \ln2/k_{\mathrm{total}}$, उसके व्युत्क्रमानुपाती।
\textbf{12.} $k_{\mathrm{total}}(\mathrm H) = 0.75\,k_{\mathrm{total}}(\mathrm H) + 0.25\,k_{\mathrm{total}}(\mathrm H)$, विमेथिलन और अन्य मार्ग।
\textbf{13.} 6.5 से भाग।
\textbf{14.} $0.25 + 0.75/6.52 = 0.365$।
\textbf{15.} $5.0/0.365 = \qty{13.7}{h}$।
\textbf{16.} 2.7।
\textbf{17.} 6.5, पूरा समस्थानिक प्रभाव।
\textbf{18.} छोटी या कम बार की खुराकें, और रक्त में अधिक स्थिर स्तर।
\textbf{19.} $\ln(k/T)$, $1/T$ के विरुद्ध।
\textbf{20.} $\Delta^{\ddagger}H^\circ = 61.9 \pm \qty{0.3}{kJ/mol}$, $\Delta^{\ddagger}S^\circ = -60.1 \pm \qty{0.9}{J.K^{-1}.mol^{-1}}$।
\textbf{21.} $\Delta^{\ddagger}H^\circ = 66.8 \pm \qty{0.2}{kJ/mol}$, $\Delta^{\ddagger}S^\circ = -60.3 \pm \qty{0.7}{J.K^{-1}.mol^{-1}}$।
\textbf{22.} $\Delta\Delta^{\ddagger}H^\circ = 4.88 \pm \qty{0.33}{kJ/mol}$ ($\sqrt{0.26^2 + 0.21^2}$)।
\textbf{23.} हाँ: 4.83 एक मानक त्रुटि के भीतर भली-भाँति आता है। समस्थानिक प्रभाव
पूर्णतः \omterm{def:b3:quantum-model-systems:oscillator}{शून्य-बिंदु ऊर्जाओं} से समझाया जाता है, \omterm{def:b3:quantum-model-systems:tunnelling}{सुरंगन} का कोई संकेत नहीं।
\textbf{24.} वे अपनी त्रुटियों के भीतर बराबर हैं: संक्रमण अवस्था की संरचना
दोनों \omterm{def:b3:rovibrational-spectroscopy:rotational-constant}{समस्थानिकरूपियों} के लिए समान है, जैसा भाग I में माना गया।
\textbf{25.} \textbf{\qty{37}{\celsius} पर $t_{1/2}(\mathrm D)/t_{1/2}(\mathrm H) \approx 2.7$}: \qty{13.7}{h}, \qty{5.0}{h} के विरुद्ध।
\end{solution}
